\chapter{Описание разработанного программного обеспечения} \label{ch2}

\section{Структура проекта} \label{ch2:structure}

Программное средство размещено в каталоге \texttt{benchmark\_app}. Проект сделан компактным: он не требует базы данных, сервера или браузера, а все измерения выполняются локально в Node.js.

\noindent
\begingroup
\centering
\small
\begin{longtable}[c]{|p{5cm}|p{10cm}|}
    \caption{Структура проекта}
    \label{tab:project-structure}
    \\
    \hline
    Файл или каталог & Назначение \\ \hline
    \endfirsthead
    \hline
    Файл или каталог & Назначение \\ \hline
    \endhead
    \hline
    \texttt{package.json} & npm-скрипты и зависимости Jest, Vitest, Mocha, pidusage. \\ \hline
    \texttt{requirements.txt} & Краткие требования к окружению и команды запуска. \\ \hline
    \texttt{src/data-loader.js} & Импорт данных из JSON и CSV, детерминированное расширение JSON-массива до заданного размера. \\ \hline
    \texttt{src/sample-analysis.js} & Чистая функция анализа числовой выборки. \\ \hline
    \texttt{src/benchmark.js} & Запуск фреймворков, сэмплирование PID, сохранение отчетов. \\ \hline
    \texttt{tests/jest} & Модульный тест для Jest. \\ \hline
    \texttt{tests/vitest} & Модульный тест для Vitest. \\ \hline
    \texttt{tests/mocha} & Модульный тест для Mocha. \\ \hline
    \texttt{data} & Эталонные и авторские входные файлы. \\ \hline
    \texttt{reports} & Автоматически сформированные результаты эксперимента. \\ \hline
\end{longtable}
\normalsize
\endgroup

\section{Импорт данных JSON и CSV} \label{ch2:input}

Проект поддерживает два формата входа. JSON-файл содержит имя набора, требуемый размер и начальные значения:
\begin{lstlisting}[breaklines=true]
{
  "name": "author-json-medium",
  "size": 5000,
  "values": [42, 7, 19, 7, 3, 11, 29, 5, 101, 17]
}
\end{lstlisting}

Если поле \texttt{size} больше длины массива \texttt{values}, недостающие элементы создаются детерминированно. Это позволяет проверить входы переменного размера без хранения больших файлов.

CSV-файл содержит колонку \texttt{value}:
\begin{lstlisting}[breaklines=true]
value
42
7
19
\end{lstlisting}

JSON используется как основной формат масштабируемого эксперимента, а CSV добавлен как альтернативный импорт для проверки требования задания.

\section{Эталонные и авторские данные} \label{ch2:data}

В проекте подготовлены три вида данных:
\begin{enumerate}
    \item \texttt{reference-input.json} --- малый эталонный вход;
    \item \texttt{reference-output.json} --- ожидаемый результат для эталонного входа;
    \item \texttt{author-example.json} и \texttt{author-example.csv} --- авторские входные файлы для запуска эксперимента.
\end{enumerate}

Эталонный файл нужен для быстрой проверки корректности вычислительной части. Авторские файлы нужны для основной демонстрации сравнения фреймворков.

\section{Средства и методы тестирования} \label{ch2:testing}

Во всех трех фреймворках тестируется один и тот же модуль \texttt{sample-analysis.js}. Тесты устроены одинаково:
\begin{enumerate}
    \item прочитать путь к входному файлу из переменной окружения \texttt{TEST\_INPUT\_FILE};
    \item загрузить данные через общий модуль \texttt{data-loader.js};
    \item выполнить анализ выборки несколько раз;
    \item проверить базовые инварианты результата.
\end{enumerate}

Повтор анализа задается переменной \texttt{TEST\_REPEAT\_COUNT}. В бенчмарке установлено значение \texttt{25}. Это сделано для того, чтобы тест выполнялся не мгновенно, а с измеримой нагрузкой на CPU.

В Jest и Vitest используется функция \texttt{expect}. В Mocha используется стандартный модуль \texttt{node:assert/strict}. Такое решение снижает число дополнительных зависимостей и делает сравнение ближе к сравнению раннеров, а не библиотек утверждений.

\section{Инструкция по установке и запуску} \label{ch2:run}

Установка зависимостей:
\begin{lstlisting}[breaklines=true]
cd benchmark_app
npm install
\end{lstlisting}

Проверка эталонного файла:
\begin{lstlisting}[breaklines=true]
npm run verify
\end{lstlisting}

Запуск отдельных фреймворков:
\begin{lstlisting}[breaklines=true]
npm run test:jest
npm run test:vitest
npm run test:mocha
\end{lstlisting}

Запуск полного бенчмарка на JSON:
\begin{lstlisting}[breaklines=true]
npm run benchmark
\end{lstlisting}

Запуск полного бенчмарка на CSV:
\begin{lstlisting}[breaklines=true]
npm run benchmark:csv
\end{lstlisting}

После запуска формируются файлы \texttt{reports/benchmark-results-json.json}, \texttt{reports/benchmark-results-json.csv}, \texttt{reports/benchmark-results-csv.json} и \texttt{reports/benchmark-results-csv.csv}.
